% Generated draft: requires mathematical and source review.
\documentclass[11pt]{article}
\usepackage[margin=25mm]{geometry}
\usepackage{fontspec}
\setmainfont{Latin Modern Roman}
\usepackage{amsmath,amssymb,mathtools,unicode-math}
\setmathfont{Latin Modern Math}
\usepackage{enumitem,needspace,xurl,hyperref}
\setlength{\parindent}{0pt}
\setlength{\parskip}{5pt}
\setlength{\emergencystretch}{3em}
\newcommand{\N}{\mathbb{N}}
\newcommand{\Z}{\mathbb{Z}}
\newcommand{\Q}{\mathbb{Q}}
\newcommand{\R}{\mathbb{R}}
\newcommand{\C}{\mathbb{C}}
\newcommand{\F}{\mathbb{F}}
\newcommand{\A}{\mathbb{A}}
\newcommand{\PP}{\mathbb{P}}
\newcommand{\E}{\mathbb{E}}
\newcommand{\eps}{\varepsilon}
\newcommand{\dd}{\,\mathrm{d}}
\begin{document}
\section*{UnsolvedMath ID: 700471}
\textbf{Source problem code:} catalogue entry
\subsection{Definitions and mathematical statement}
A group identity is an equation \(w=1\), where \(w\) is a word in finitely many variables and their inverses. A group satisfies this identity if every substitution of its elements for the variables gives the identity element. A variety of groups is the class of all groups satisfying some specified set of group identities. A subvariety of a variety \(\mathcal V\) is a variety contained in \(\mathcal V\). Write \(\operatorname{Sub}(\mathcal V)\) for its set of subvarieties, ordered by inclusion.

For a group \(G\), put \(G^{(0)}=G\) and \(G^{(n+1)}=[G^{(n)},G^{(n)}]\). A soluble variety is a variety \(\mathcal V\) for which there exists an integer \(d\geq 0\) such that \(G^{(d)}=1\) for every \(G\in\mathcal V\). No particular derived-length bound is supplied.

Consider the following properties of \(\mathcal V\):
\begin{enumerate}
\item Every descending sequence of subvarieties
\[
\mathcal V_0\supseteq\mathcal V_1\supseteq\mathcal V_2\supseteq\cdots,
\qquad \mathcal V_n\subseteq\mathcal V,
\]
is eventually constant: there exists \(N\geq 0\) such that \(\mathcal V_n=\mathcal V_N\) for every \(n\geq N\). This is the minimum condition for subvarieties.
\item \(\lvert\operatorname{Sub}(\mathcal V)\rvert\leq\aleph_0\), including the possibility of finitely many subvarieties.
\item \(\mathcal V\) has no infinite independent system of identities.
\end{enumerate}

The supplied wording does not specify the meaning of ``has'' or the ambient class used for independence in property~(3). The following two precise readings must therefore be distinguished. For a set \(\Sigma\) of identities, define
\[
\mathcal V[\Sigma]=\{G\in\mathcal V: G\text{ satisfies every identity in }\Sigma\}.
\]
A system \(\Sigma\) is independent relative to \(\mathcal V\) if, for every \(\sigma\in\Sigma\), there exists a group \(G_\sigma\in\mathcal V\) satisfying every identity in \(\Sigma\setminus\{\sigma\}\) and failing \(\sigma\). Equivalently,
\[
\mathcal V[\Sigma]\subsetneq\mathcal V[\Sigma\setminus\{\sigma\}]
\qquad(\sigma\in\Sigma).
\]
Under the relative reading, property~(3) asserts that no infinite system is independent relative to \(\mathcal V\). Another possible reading asserts that there is no infinite set \(\Sigma\) of identities all valid in \(\mathcal V\) that is independent over the class of all groups; here each witness \(G_\sigma\) may lie outside \(\mathcal V\). These readings are not identified by the supplied record.

\textbf{Question (Yu. G. Kleiman, Kourovka Notebook 8.19).} Which pairs among properties~(1), (2), and~(3) are equivalent as properties of varieties of groups? The source writes ``(soluble) varieties'', leaving the role of that parenthetical restriction unresolved. Thus the equivalence question should be retained both for arbitrary group varieties and for soluble group varieties, pending clarification of the intended scope and of property~(3). For any fixed scope and reading, equivalence of properties \(i\) and \(j\) means that \(P_i(\mathcal V)\Longleftrightarrow P_j(\mathcal V)\) for every variety \(\mathcal V\) in that scope. No equivalence or implication is asserted here.
\subsection{Short English statement}
Kleiman asks which of three conditions on a group variety are equivalent: the descending chain condition on its subvarieties, having at most countably many subvarieties, and having no infinite independent system of identities. The source leaves both the parenthetical soluble restriction and the precise interpretation of independence unclear; these qualifications require review.
\subsection{Sources}
\begin{itemize}
\item Supplied local UnsolvedMath registry record, permanent ID \(700471\), exact source problem code \texttt{catalogue entry}, title \emph{Kourovka Notebook Problem 8.19}. The registry labels the problem \emph{open}; this is supplied source metadata, not an independently verified mathematical or literature assessment.
\item \emph{The Kourovka Notebook (Unsolved Problems in Group Theory)}, original problem 8.19, 8th issue (1982); proposed by Yu. G. Kleiman.
\item The supplied background identifies the 21st issue edition (2026), \texttt{arXiv:1401.0300v48}. The supplied collection description names E. I. Khukhro and V. D. Mazurov as editors. This bibliographic attribution is reproduced from the record; the cited edition has not been independently consulted.
\end{itemize}
\end{document}
